% Editable TikZ for all A5 worked examples. Curves with numerical choices say so.
\newcommand{\QAOFrame}[3]{\begin{qualityfigure}#1\bookcaption{#2}\label{#3}\end{qualityfigure}}
\pgfplotsset{qualityoptics/.style={axis lines=left,grid=major,width=.74\linewidth,height=5cm,
 tick label style={font=\small},label style={font=\small},every axis plot/.append style={thick}}}
\newcommand{\QAOFermat}[1]{\begin{tikzpicture}[>=Latex,scale=.85]
 \foreach\y in {-.5,.5,1.5}{\draw[guide](-3,\y)--(3,\y);}
 \coordinate(A)at(-2.7,1.1);\coordinate(B)at(2.7,1.1);\coordinate(P)at(0,0);
 \draw[DeepBlue,thick](A)..controls(-1,-.3666667)and(1,-.3666667)..(B);
 \fill(A)circle(2pt)node[left]{$A$};\fill(B)circle(2pt)node[right]{$B$};
 \node[above]at(0,1.8){$n=n(\vect r)>0$};
 \ifnum#1=1
  \draw[Gold!80!black,dashed](A)..controls(-1,.2333333)and(1,.2333333)..(B);
  \draw[vector](P)--(1.1,0)node[below]{$\vect\tau$};\draw[vector](P)--(0,.45)node[right]{$\delta\vect r$};
  \node[below]at(0,-1.15){固定端点：$\delta\vect r(A)=\delta\vect r(B)=0$};
 \else\node[below]at(0,-1.15){端点固定，在所有光滑连接路径中考察光程驻值};\fi
\end{tikzpicture}}
\newcommand{\QAOBeachAnswer}{\begin{tikzpicture}[>=Latex,scale=.85]
 \fill[PaleBlue](-.4,0)rectangle(6.7,2.6);\draw[axis](-.4,0)--(7,0)node[right]{$x$};\draw[axis](0,0)--(0,2.9)node[above]{$y$};
 \draw[DeepBlue,very thick,->,domain=90:179.5,samples=90]plot({2*(\x*pi/180-.5*sin(2*\x))},{2*sin(\x)^2});
 \draw[guide](0,2)--(3.14159,2);\node[left]at(0,2){$y_m$};\node[below]at(6,0){海岸};
 \node at(1.5,.7){$H(y)=ay$};\node[above]at(4.3,2.3){向岸分支};
\end{tikzpicture}}
\newcommand{\QAOBeachQuestion}{\begin{tikzpicture}[>=Latex,scale=.85]
 \fill[PaleBlue](-.5,0)rectangle(5.8,2.6);\draw[axis](-.5,0)--(6.2,0)node[right]{$x$，海岸};\draw[axis](0,0)--(0,3)node[above]{$y$};
 \foreach\y in {.7,1.4,2.1}{\draw[guide](0,\y)--(5.6,\y);}
 \fill(3.2,1.8)circle(2pt)node[above]{$P$};\draw[vector](3.2,1.8)--(4.1,1.1)node[right]{波线方向};
 \draw[guide](3.2,1.8)--(3.2,.5);\node at(3.45,1.25){$\theta$};\node[above]at(1.4,2.4){$H(y)=ay$};
\end{tikzpicture}}
\newcommand{\QAOFishQuestion}{\begin{tikzpicture}[>=Latex,scale=.85]
 \draw[guide](0,0)circle(1.5);\draw[guide](0,0)circle(2.2);\fill(0,0)circle(1.5pt)node[below]{$O$};
 \fill(.7,.3)circle(2pt)node[above right]{$P\ne O$};\draw[vector](.7,.3)--(1.5,.9);\draw[vector](.7,.3)--(.2,1);\draw[vector](.7,.3)--(-.2,.1);
 \draw[->](0,0)--(2.03,-.85)node[midway,below]{$r$};
 \node[anchor=west,align=left]at(3,.5){待构造 $n(r)$\\不同初始方向的光线\\寻找共同像点 $P'$};
\end{tikzpicture}}
\newcommand{\QAOFishAnswer}{\begin{tikzpicture}[>=Latex,scale=.85]
 \draw[guide](0,0)circle(1);\draw[axis](-2.8,0)--(2.1,0)node[right]{$x/a$};
 \foreach\h in {-.9,0,.9}{\pgfmathsetmacro{\rr}{sqrt(1.5625+\h*\h)}\draw[DeepBlue](-.75,\h)circle(\rr);}
 \fill(.5,0)circle(2pt)node[above right]{$P$};\fill(-2,0)circle(2pt)node[above left]{$P'$};\fill(0,0)circle(1pt)node[below]{$O$};
 \node[anchor=west,align=left]at(2.6,.4){示例 $z_0/a=1/2$\\$z_i/a=-2$\\虚线：$r=a$\\实线：球面大圆的投影};
\end{tikzpicture}}
\newcommand{\QAOHelixAnswer}{\begin{tikzpicture}\begin{axis}[qualityoptics,xmin=0,xmax=2,ymin=0,ymax=2.5,
 xlabel={$s=z/f$},ylabel={$\rho/r$},legend style={font=\small,at={(.03,.97)},anchor=north west}]
 \addplot[DeepBlue,domain=0:2,samples=80]{sqrt(1+x*x)};\addlegendentry{不加透镜}
 \addplot[Gold!80!black,domain=0:2,samples=80]{sqrt((1-x)*(1-x)+x*x)};\addlegendentry{加薄透镜}
 \addplot[only marks,mark=*]coordinates{(.5,.7071068)};
 \draw[dashed](axis cs:.5,0)--(axis cs:.5,.7071068);\node[above right]at(axis cs:.5,.7071068){束腰};
\end{axis}\end{tikzpicture}}
\newcommand{\QAOThickQuestion}{\begin{tikzpicture}[>=Latex,x=.95cm,y=1cm,font=\small]
 \draw[axis](-1,0)--(9,0);\draw[very thick](0,-1)--(0,1) (3,-1)--(3,1) (8,-1.2)--(8,1.2);
 \fill[PaleBlue](.05,-.95)rectangle(2.95,.95);\node at(1.5,.35){$n=3$};
 \node[above]at(0,1.1){左面 $R$};\node[above]at(3,1.1){右面 $R'$，刻线};\node[above]at(8,1.2){物镜 $L_O$};
 \draw[<->](0,-1.4)--(3,-1.4)node[midway,below]{$x$};\draw[<->](3,-1.4)--(8,-1.4)node[midway,below]{$y$};
 \draw[<->](0,-2)--(8,-2)node[midway,below]{$L=8\,\mathrm{cm}$};
 \draw[vector](7,1.8)--(4.6,1.8)node[midway,above]{传播方向};\node[left]at(-.3,.3){眼瞳 $d$};
\end{tikzpicture}}
\newcommand{\QAOPeriscopeQuestion}{\begin{tikzpicture}[>=Latex,scale=.75,font=\small]
 \draw[guide](-2,4)--(0,4)--(0,-1)--(2,-1);
 \draw[very thick](-.45,4.45)--(.45,3.55) (-.45,-.55)--(.45,-1.45);
 \foreach\y/\i in {3.4/0,2.5/1,1.3/2,.1/3}{\draw[very thick](-.65,\y)--(.65,\y);\node[right]at(.8,\y){$L_\i$};}
 \draw[very thick](1.6,-1.8)--(1.6,-.2)node[right]{$L_4$};
 \draw[<->](-2.2,-1)--(-2.2,4)node[midway,left]{$H$};\node[right]at(.7,4){$P$};\node[right]at(.7,-1){$P'$};
 \node[anchor=west,align=left]at(3,1.8){$L_1$：第一实像面 $y'$\\$L_3$：第二实像面 $y''$\\各镜的 $f_i,D_i$ 待配合设计};
\end{tikzpicture}}
\newcommand{\QAOInterfaces}[1]{\begin{tikzpicture}[>=Latex,scale=.85,font=\small]
 \fill[PaleBlue](0,-1.7)rectangle(2,1.7);\draw[very thick](0,-1.7)--(0,1.7) (2,-1.7)--(2,1.7);
 \node[above]at(0,1.7){界面1};\node[above]at(2,1.7){界面2};\draw[<->](0,-2)--(2,-2)node[midway,below]{单程相位 $\delta$};
 \draw[vector](-2,1.2)--(0,.8);\draw[vector](0,.8)--(-1.8,.4)node[left]{$r_1$};
 \draw[vector](0,.8)--(2,.4);\draw[vector](2,.4)--(3.5,.1)node[right]{透射};
 \ifnum#1=1
  \draw[vector,Gold!80!black](2,.4)--(0,0)--(2,-.4)--(3.5,-.7);
  \draw[vector,Gold!80!black](0,0)--(-1.8,-.4)node[left]{$t_1t'_1r_2\e^{2\ii\delta}$};
  \node[below,align=center]at(0,-2.8){每增加一次往返，振幅乘 $r'_1r_2\e^{2\ii\delta}$};
 \else\node[below,align=center]at(0,-2.8){系数均按传播通道的能流归一化；画线只示传播顺序};\fi
\end{tikzpicture}}
\newcommand{\QAOStackQuestion}{\begin{tikzpicture}[>=Latex,scale=.85]
 \foreach\i in {0,1,2}{\fill[PaleBlue](2*\i,-1)rectangle(2*\i+.8,1);\fill[SoftGray](2*\i+.8,-1)rectangle(2*\i+2,1);\node at(2*\i+.4,0){$n_1$};\node at(2*\i+1.4,0){$n_2$};}
 \draw[<->](0,-1.3)--(.8,-1.3)node[midway,below]{$d_1$};\draw[<->](.8,-1.3)--(2,-1.3)node[midway,below]{$d_2$};
 \draw[vector](-1.8,0)--(-.2,0);\draw[axis](0,1.5)--(6.5,1.5)node[right]{$z$};\node at(6.7,0){$\cdots$};
\end{tikzpicture}}
\newcommand{\QAOStackAnswer}{\begin{tikzpicture}\begin{axis}[qualityoptics,xmin=0,xmax=2,ymin=-1.5,ymax=1.3,
 xlabel={$\omega/\omega_0$},ylabel={$(\operatorname{tr}M)/2$}]
 \addplot[draw=none,fill=Gold!18]coordinates{(.78365,-1.5)(1.21635,-1.5)(1.21635,1.3)(.78365,1.3)}\closedcycle;
 \addplot[DeepBlue,domain=0:2,samples=100]{1-2.25*sin(90*x)^2};
 \addplot[dashed,domain=0:2]{-1};\addplot[dashed,domain=0:2]{1};\node at(axis cs:1,.6){禁带};
\end{axis}\end{tikzpicture}}
\newcommand{\QAOCascadeAnswer}{\begin{tikzpicture}
 \begin{axis}[qualityoptics,width=.44\linewidth,height=4.5cm,xmin=-3,xmax=3,ymin=0,ymax=1.1,
 xlabel={$\sqrt{2019}(\pi\Delta\ell/\lambda-m\pi)$},ylabel={$I_O/I_{\max}$}]
  % Double-precision samples of cos(x/sqrt(2019))^4038 avoid amplifying
  % TeX's trigonometric rounding into false ripples at the spectral peak.
  \addplot[DeepBlue,smooth]coordinates{
   (-3,.0001226)(-2.9,.0002213)(-2.8,.0003917)(-2.7,.0006793)(-2.6,.0011549)
   (-2.5,.0019242)(-2.4,.0031425)(-2.3,.0050301)(-2.2,.0078918)(-2.1,.0121357)
   (-2,.0182915)(-1.9,.0270227)(-1.8,.0391300)(-1.7,.0555379)(-1.6,.0772629)
   (-1.5,.1053552)(-1.4,.1408137)(-1.3,.1844760)(-1.2,.2368872)(-1.1,.2981612)
   (-1,.3678491)(-.9,.4448340)(-.8,.5272746)(-.7,.6126143)(-.6,.6976689)
   (-.5,.7787968)(-.4,.8521420)(-.3,.9139306)(-.2,.9607893)(-.1,.9900498)
   (0,1)(.1,.9900498)(.2,.9607893)(.3,.9139306)(.4,.8521420)(.5,.7787968)
   (.6,.6976689)(.7,.6126143)(.8,.5272746)(.9,.4448340)(1,.3678491)
   (1.1,.2981612)(1.2,.2368872)(1.3,.1844760)(1.4,.1408137)(1.5,.1053552)
   (1.6,.0772629)(1.7,.0555379)(1.8,.0391300)(1.9,.0270227)(2,.0182915)
   (2.1,.0121357)(2.2,.0078918)(2.3,.0050301)(2.4,.0031425)(2.5,.0019242)
   (2.6,.0011549)(2.7,.0006793)(2.8,.0003917)(2.9,.0002213)(3,.0001226)};
  \addplot[dashed,domain=-3:3]{.5};
 \end{axis}
 \begin{axis}[qualityoptics,at={(8cm,0)},anchor=south west,width=.44\linewidth,height=4.5cm,
 xmin=-2,xmax=2,ymin=0,ymax=1.1,xlabel={$Dy/(\lambda L_1)$},ylabel={$I(y)/I(0)$}]
  \addplot[Gold!80!black,domain=-2:2,samples=110]{cos(180*x)^2};
 \end{axis}
\end{tikzpicture}}
\newcommand{\QAOBiprismAnswer}{\begin{tikzpicture}[>=Latex,x=1.2cm,y=1cm]
 \draw[axis](-.4,0)--(4.6,0)node[right]{$y$};\draw[axis](0,-.1)--(0,1.9)node[above]{$\Phi-C$};
 \foreach\i in {0,1,2,3}{\draw[DeepBlue,thick,domain=0:1,samples=25]plot({\i+\x},{4*(\x-.5)*(\x-.5)});}
 \foreach\i in {1,2,3}{\draw[Gold!80!black,dashed](\i,0)--(\i,1.3);\node[above]at(\i,1.3){$b_\i$};}
 \draw[<->](.5,-.5)--(1.5,-.5)node[midway,below]{$d$};\node[below]at(2,-1.1){四单元示意：三条棱脊产生四段相同的平移相位};
\end{tikzpicture}}
\newcommand{\QAOGratingAnswer}{\begin{tikzpicture}[>=Latex,scale=.85,font=\small]
 \begin{scope}
  \draw[axis](-.3,0)--(4.5,0)node[right]{$\Re A$};\draw[axis](0,-.3)--(0,2)node[above]{$\Im A$};
  \foreach\i in {0,...,7}{\draw[vector](.5*\i,0)--(.5*\i+.5,0);}
  \node[below]at(1.8,-.7){主峰：$\xi=0$，八个同相振幅};
 \end{scope}
 \begin{scope}[xshift=8cm]
  \draw[axis](-.8,0)--(1.9,0)node[right]{$\Re A$};\draw[axis](0,-.3)--(0,2)node[above]{$\Im A$};
  \foreach\a/\b/\c/\d in {0/0/.65/0,.65/0/1.1096/.4596,1.1096/.4596/1.1096/1.1096,1.1096/1.1096/.65/1.5692,.65/1.5692/0/1.5692,0/1.5692/-.4596/1.1096,-.4596/1.1096/-.4596/.4596,-.4596/.4596/0/0}{\draw[vector](\a,\b)--(\c,\d);}
  \node[below]at(.4,-.7){第一零点：$\xi=2\pi/8$，向量闭合};
 \end{scope}
\end{tikzpicture}}
\newcommand{\QAOMichelsonQuestion}{\begin{tikzpicture}[>=Latex,scale=.85]
 \draw[very thick](0,-1.2)--(0,1.2) (1,-1.2)--(1,1.2);\draw[<->](0,-1.6)--(1,-1.6)node[midway,below]{$d$};
 \draw[vector](-2.5,1.2)--(0,.3);\draw[vector](0,.3)--(-2.5,-.6);
 \draw[vector,dashed](0,.3)--(1,-.06);\draw[vector](1,-.06)--(-1.5,-.96);
 \node[above]at(.5,1.8){迈克尔逊的两等效反射面};
 \node[anchor=west,align=left]at(2,.4){同一望远镜视场\\初态：20亮环\\吞入20环后：10亮环\\中心亮斑另计};
\end{tikzpicture}}
\newcommand{\QAOCylindersQuestion}{\begin{tikzpicture}[>=Latex,scale=.85]
 \begin{scope}[xshift=-3cm]
  \draw[axis](-1.8,0)--(2,0)node[right]{$x$};\draw[axis](0,-.5)--(0,2)node[above]{$z$};
  \draw[thick,domain=-1.5:1.5]plot(\x,{.5*\x*\x});\node[above left]at(-1.3,1.1){上：$R_2$};
  \draw[vector](1,2)--(1,1.3);\node[below]at(0,-.7){下圆柱沿 $x$ 方向平直};
 \end{scope}
 \begin{scope}[xshift=3.8cm]
  \draw[axis](-1.8,0)--(2,0)node[right]{$y$};\draw[axis](0,-1.8)--(0,1.3)node[above]{$z$};
  \draw[thick,domain=-1.5:1.5]plot(\x,{-.5*\x*\x});\node[below left]at(-1.3,-1){下：$R_1$};
  \draw[vector](1,1.2)--(1,.4);\node[below]at(0,-2){上圆柱沿 $y$ 方向平直};
 \end{scope}
\end{tikzpicture}}
